% ==================================================
% CHAPTER 1: INTRODUCTION
% ==================================================

\chapter{Introduction}
\label{ch:introduction}

\section{Background}

At Thapar Institute every hostel resident is allotted one room for the
semester, and the allotment is fixed once it is made. A student who ends
up in the wrong hostel, on the wrong floor, facing the wrong direction,
in the wrong room type or away from their friends has no structured way
to change it. The informal remedy is to ask around or post in a group
chat. That only ever surfaces a handful of offers, and a workable trade
almost always needs more than two people.

A direct one-to-one swap needs two students who each want exactly what
the other has. That coincidence is rare. What usually does exist is a
longer exchange:

\begin{center}
\textit{A takes B's room, B takes C's room, and C takes A's room.}
\end{center}

Every member of that rotation ends up better off, but nobody can see it
from inside their own conversation, because seeing it needs every other
student's room and every other student's preferences at once. This
project treats room swapping as a \emph{visibility} problem rather than
a negotiation problem: the service holds all rooms and all preferences
centrally, computes the exchange cycles hidden in that data, and offers
them back to the students as concrete swaps they can approve.

\section{Problem Statement}

A hostel room exchange service has to solve several problems at once:

\begin{enumerate}
  \item \textbf{Express what students want.} A student cares about the
  hostel, floor, direction, washroom type, room type and roommate, some
  as absolute requirements and some as preferences of different weight.
  \item \textbf{Find exchanges nobody can see alone.} The service must
  consider every participant's room and preferences together and find
  multi-person chains, not only pairs.
  \item \textbf{Never make anyone worse off.} A student will not approve
  a swap that leaves them in a worse room, so an offered chain must
  improve or keep every member's position, as a property of the system
  rather than an accident.
  \item \textbf{Offer chains that can actually happen.} Every member of a
  chain has to agree before anyone moves. A 28-person chain is one offer
  with 28 chances to fail, so short chains matter as much as good ones.
  \item \textbf{Handle people who refuse or never answer.} One refusal or
  one silent student must not block everyone else indefinitely.
  \item \textbf{Keep an auditable record} of who held which room and
  when, which no later bug can rewrite.
\end{enumerate}

\section{Objectives}

The objectives set in the proposal, and refined after the mid-semester
evaluation, were:

\begin{enumerate}
  \item Model room swapping as a closed-pool assignment problem, so that
  swap chains fall out of the assignment's cycle structure instead of
  being searched for separately.
  \item Solve for the best assignment with a proven combinatorial
  algorithm, and generate alternative assignments so that shorter, more
  executable chains can be offered.
  \item Filter every candidate chain so that nobody in it is worse off.
  \item Give every student their own best options, which differ in kind
  and not only in score.
  \item Model rounds and proposals as explicit state machines, including
  rejection and expiry.
  \item Persist the data in a managed relational database, with the
  important invariants enforced by the database itself.
  \item Provide web interfaces for students and wardens.
  \item Evaluate the engine with reproducible cohorts and metrics that
  describe the whole cohort, and use the results to set its parameters.
  \item Follow sound software-engineering practice: version control,
  automated tests, continuous integration, schema migrations and
  documentation.
\end{enumerate}

\section{From Prototype to Final Product}

The mid-semester report described a working prototype: the matching
engine, the round and proposal state machines, a FastAPI backend and
plain HTML pages. The final product keeps the engine and rebuilds or
extends most of the rest. Table~\ref{tab:prototype-vs-final} summarises
the changes.

\begin{table}[H]
\centering
\small
\begin{tabularx}{\textwidth}{>{\raggedright\arraybackslash}p{0.22\textwidth} L L}
\toprule
\textbf{Aspect} & \textbf{Prototype (mid-semester)} & \textbf{Final product} \\
\midrule
What is evaluated & One best plan for the whole hostel & Also the chains offered to each student: up to three each, the best available to them \\
Chain length & Unlimited & Capped at eight students, chosen by a 300-student benchmark \\
Proposal expiry & In the state model only & Enforced: late replies refused, overdue proposals swept to expired \\
Locked-out members & Members of settled chains stayed blocked & Members freed when their proposal settles (migration 0002) \\
Roommates & Scored, not requestable & Mutual roommate requests, sent and accepted in the app \\
Web interface & Server-rendered HTML and JavaScript & React + TypeScript single-page app built with Vite \\
Schema & One \texttt{schema.sql} & Numbered migrations with a runner, plus seed scripts \\
Production checks & None & Secure cookies, refuses to start without a database \\
Murty's search & Full-size subproblems & Subproblems shrink as edges are fixed: 4--13$\times$ faster \\
Warden controls & Create, open, run & Also cancel a round and withdraw a ring \\
Deadlines & Not shown & Countdown per offer; overdue offers shown as expired \\
Tests & 127 & 152, four against PostgreSQL in CI \\
Evaluation & One 30-student cohort & Plus ten seeded 300-student cohorts per chain cap, and a scale study \\
\bottomrule
\end{tabularx}
\caption{The prototype and the final product}
\label{tab:prototype-vs-final}
\end{table}

\section{Organisation of the Report}

Chapter~\ref{ch:requirements} sets out the requirements and the
development process. Chapter~\ref{ch:design} presents the architecture,
data flow, data model and state machines, and
Chapter~\ref{ch:implementation} the implementation of the engine, the
round workflow and the interfaces. Chapter~\ref{ch:security} covers
security, Chapter~\ref{ch:testing} testing and the defects it found, and
Chapter~\ref{ch:evaluation} the evaluation and the parameters it set.
Chapter~\ref{ch:deployment} describes configuration and deployment, and
Chapter~\ref{ch:conclusion} concludes with limitations and future work.
